\section{File \texttt{NS\_conversion.txt}}
\subsection{Purpose}
This is text a file specifying, for each nodestring, the name of the boundary and/or qsection in which it will be converted.

\subsection{Format description}
\begin{footnotesize}
\begin{tabular}{p{4cm} | p{10cm} }
\hline
Field(s) & Specification \\
\hline
\hline
$NS$    $bc\_code$   $qsection\_name$ & This is the header line. It must be inserted without modification \\
\hline
$ns_1 \rightharpoonup bc_1 \rightharpoonup qsection_1$ &
\begin{itemize}
\item [$ns_1$] SMS ID of the 1st nodestring
\item [$bc_1$] BC code of the destination boundary for the 1st nodestring, $0$ means no boundary
\item [$qsection_1$] Name of the destination qsection for the 1st nodestring, $None$ means no qsection
\item [$\rightharpoonup$] Tabulation character or space
\end{itemize} \\
\hline
$\vdots$ & (repeat) \\
\hline
$ns_N \rightharpoonup bc_N \rightharpoonup qsection_N$ &
\begin{itemize}
\item [$ns_N$] SMS ID of the $N$th nodestring
\item [$bc_N$] BC code of the destination boundary for the $N$st nodestring, $0$ means no boundary
\item [$qsection_N$] Name of the destination qsection for the $N$st nodestring, $None$ means no qsection
\end{itemize} \\
\hline
\end{tabular}
\end{footnotesize}

\subsection{Sample file}

In the following file : \begin{itemize}
    \item Nodestring 1 is converted into a boundary with BC code $-1$ and a qsection with name $q_1$.
    \item Nodestring 2 is converted into a qsection with name $q_2$.
    \item Nodestring 3 is converted into a boundary with BC code $-2$.
\end{itemize}

\begin{verbatim}
NS	bc_code	qsection_name
1	-1	q_1
2	0	q_2
3	-2	None
\end{verbatim}
%%